\documentclass{article}
\usepackage{tikz}
\usetikzlibrary{angles, quotes}
\usepackage{amssymb}
\usepackage{pgfplots}
\pgfplotsset{width=10cm,compat=1.9}
\usepgfplotslibrary{statistics}
\usepackage[left=2.54cm,right=2.54cm,top=2.54cm,bottom=2.54cm]{geometry}
\usepackage{tasks} %Paquete necesario para la producción de items horizontales para algunos ejercicios de tarea empleando el comando /begin{multicols}}{#}.../end{multicols} para ayudarnos a reducir las páginas
\usepackage{pdflscape} %comando que permite cambiar a orientación horizontal las hojas%
\usepackage{soul} %Paquete para permitir la compilación de acentos usando el comando {\hl{}}}\\
\sethlcolor{yellow(munsell)} %comando para definir el color para el subrayado usando el comando \hl
\usepackage[utf8]{inputenc}
\usepackage[spanish]{babel}
\usepackage{icomma}
\usepackage{ragged2e}
\usepackage{siunitx}
\usepackage{url}
\usepackage[colorlinks=true, urlcolor=blue,  linkcolor=black, citecolor=green]{hyperref}
\usepackage{pdfpages}
\usepackage{blindtext} %Paquete que genera texto ficticio (dummy text) con el comando: \blindtext
\usepackage{cancel} %in the preamble gives you four different modes of striking through: \cancel{text to cancel} draws a diagonal line (slash) through its argument, \bcancel{text to cancel} uses the negative slope (a backslash), \xcancel{text to cancel} draws an X (actually \cancel plus \bcancel), \cancelto{〈value〉}{〈expression〉} draws a diagonal arrow through the 〈expression〉pointing to the 〈value〉 (math-mode only)
\usepackage{csquotes}
\usepackage{afterpage}
\usepackage{parskip} 
\usepackage{float}
\usepackage{enumitem}
\usepackage{multicol}%Paquete que permite la creación aislada de columnas de texto. De esta forma se reduce la cantidad de páginas en nuestro documento
\usepackage{lipsum} %Paquete que genera texto de relleno al igual que \usepackage{blindtext}, se usa con el comando \lipsum[#-#] los #(hashtags) delimitan el número de páginas que deseamos ocupar del paquete lipsum para usarlos como relleno. 

% Definir el comando personalizado para la nota
\newcommand{\mynote}[1]{%
\begin{tikzpicture}[baseline=-0.75ex]
\node[draw, circle, fill=Apple Green!60, inner sep=2pt] (note) {\textbf{Nota:}};
\end{tikzpicture}%
\ \textit{#1}%
}

\newenvironment{Figura}
{\par\medskip\noindent\minipage{\linewidth}}
{\endminipage\par\medskip}
\usepackage{caption}
\usepackage[
backend=biber,
sorting=none,
url=true
]{biblatex} %bibliografía
\addbibresource{biblio.bib}
% Establecer el espacio entre las entradas de la bibliografía
\setlength{\bibitemsep}{1\baselineskip} % Puedes ajustar el valor según tus preferencias
\usepackage{amsmath,amsthm,amssymb,amsfonts}
\usepackage{pifont} %Permite la compilación del comando \xmark: es el símbolo de la tachita
\usepackage{mathtools} %permite la compilación de símbolos para matrices
\usepackage{empheq} %Paquete que se relaciona con \usepackage[most]{tcolorbox} para la creación de cuadros/cajas de colores para delimitar los resultados matemáticos o líneas de texto usando la línea de comando: \begin{empheq}[box={\mymath[colback=orange(webcolor)!70,drop lifted shadow]}]{equation*} \end{empheq}
\usepackage[most]{tcolorbox} %Paquete que se relaciona con \usepackage{empheq} para la creación de cuadros/cajas de colores para delimitar los resultados matemáticos o líneas de texto
\renewcommand{\qedsymbol}{$\blacksquare$}
\usetikzlibrary{positioning,decorations.pathreplacing} %Paquete que permite la compilación de llaves para cuadros sinópticos (brace diagram) 
\usepackage{schemata} %The schemata package is designed just for these brace diagrams
\usepackage{booktabs}

\newcommand\AB[2]{\schema{\schemabox{#1}}{\schemabox{#2}}} %comando o paquete necesario para crear cuadros sinópticos

\newtcbox{\mymath}[1][]{%
nobeforeafter, math upper, tcbox raise base,
enhanced, colframe=blue!30!black,
colback=blue!30, boxrule=1pt,
#1} %Comando importante para encasillar los resultados matemáticos en cajas de diferentes colores, se relaciona con los paquetes \usepackage{empheq} y \usepackage[most]{tcolorbox} 

\newenvironment{sysmatrix}[1]
{\left(\begin{array}{@{}#1@{}}}
{\end{array}\right)}
\newcommand{\ro}[1]{%
\xrightarrow{\mathmakebox[\rowidth]{#1}}%
}
\newlength{\rowidth}% row operation width
\AtBeginDocument{\setlength{\rowidth}{3em}}  %Comando importante para laproducción de lineas de operaciones entre matrices, método de Gauss o Gauss Jordan se relaciona con \newenvironment{sysmatrix}

%formato para cambiar el horario a español
\usepackage[spanish]{datetime2}
\DTMsetdatestyle{spanish}

\renewcommand{\today}{\DTMdisplaydate{\the\year}{\the\month}{\the\day }{-1}}
%%%%%%%%%%%%%%%%%%%%%%%%%%%%%%%%%%%%%%%%%

\usepackage{datetime}
\newcommand{\mycurrenttime}{\xxivtime}

%%%%%%%%%%%%%%%%%%%%%%%%%%%%%%%%%%%%%%%%%%%%%%%%%%%%%%%%%%%%%%%%%%%%%%%%%%%%%%%%%%%%%%%%%%%%%%%%%%%%%%%%%%%%%%%%%%%%%%%%%%%%%%%%%%%%%%%%%%%%%%%%%


%%%%%%%%%%%%%Caja de color para el título%%%%%%%%%%%%%%%%%%%%%%%%%%%%%%%%%%%%%%%%%%%%%%%%%%%%%%

\definecolor{myframecolor}{RGB}{1, 45, 75} % Prussian Blue
\definecolor{myboxcolor}{RGB}{0, 107, 92} % Apple Green

\newtcolorbox{mybox}{
enhanced,
colback=myboxcolor!7, % Color de fondo del cuadro
colframe=myframecolor, % Color del marco del cuadro
arc=0pt,
boxrule=1pt,
borderline west={2mm}{-10mm}{myframecolor}, % Borde en el lado izquierdo
sharp corners=southwest,
width=\linewidth,
before=\par\vspace{\bigskipamount}, % Espacio antes del cuadro
after=\par\vspace{\bigskipamount} % Espacio después del cuadro
}


\newcommand{\euler}{e} %Comando para producir letra e de euler
\newenvironment{remark}{\par\vfill\footnotesize % Vertical white space above the remark and smaller font size
\begin{list}{}{
	\leftmargin=80pt % Indentation on the left
	\rightmargin=60pt}\item\ignorespaces % Indentation on the right
\makebox[-2.5pt]{\begin{tikzpicture}[overlay]
		\node[draw=Horizon!60,line width=2.5pt,circle,fill=Horizon!25,font=\sffamily\bfseries,inner sep=4pt,outer sep=0pt] at (-19pt,5pt){\textcolor{Horizon}{Nota}};\end{tikzpicture}} % Blue Nota in a circle
\advance\baselineskip -1pt}{\end{list}\vskip5pt} % Tighter line spacing and white space after remark
\usepackage{graphicx}

\usepackage{titling}

\input{colores.tex}

%%Producir signos de cita textual``Asignatura''

\newcommand{\xmark}{\ding{55}} %%%comando que genera la tachita

\newenvironment{MyColorPar}[1]{%
\leavevmode\color{#1}\ignorespaces%
}{%
}%

%%%%%%%%%%%%%%%%%%% Título de la tarea, Nombre de alumno

\title{Clase \# 29 - Estado Sólido}
\author{\emph{JA} $\boldsymbol{\mid}$}
\date{\today}

%%%%%%%%%%%%%%%%%%% Datos de la Materia

\usepackage{fancyhdr}
\fancypagestyle{plain}{%  the preset of fancyhdr 
\fancyhf{} % clear all header and footer fields
\fancyfoot[R]{
\includegraphics[width=2cm]{LogoFCFMBUAP (1).png}}
\fancyfoot[L]{{\bfseries{\thedate{}}} a las {\bfseries{\mycurrenttime{}}} horas{} (GMT-6, H. Puebla de Zaragoza, Pue)}
\fancyhead[L]{Estado sólido}
\fancyhead[R]{\theauthor}
}
\makeatletter
\def\@maketitle{%
\newpage
\null
\vskip 1em%
\begin{center}%
\let \footnote \thanks
{\LARGE \@title \par}%
\vskip 1em%
%{\large \@date}%
\end{center}%
\par
\vskip 1em}
\makeatother

\usepackage{lipsum}  
\usepackage{cmbright}


%%%%%%%%%%%%%%%%%%%%%%%%%%%%%%%%%%%%%%%%%%%%%%%%%%%%%%%%%%%%%%%%%%%%%%%%%%%%%%%%%%%%%%%%%%%%%%%%%%%%%%%%%%%%%%%%%%%%%%%%%%%%%%%%%%%%%%%%%%%%%%%%%%%%%%%%%%%%%%%%%%%%%%%%%%%%%%%%%%%%%%%%%%%%%%%%%%%%%%

\begin{document}



\maketitle


%%%%%%%%%%%%%%%%%%% Datos del profesor, campus y aula

\begin{mybox}
\noindent\begin{tabular}{@{}ll}
	{\bfseries{Profesor}} & Dra. Ana Lilia González Ronquillo\\
	\textcolor{prussianblue}{\bfseries{Campus}} / \textcolor{trueblue}{\bfseries{Facultad}}  & \textcolor{prussianblue}{\bfseries{C.U. BUAP}} /  \textcolor{trueblue}{\bfseries{FCFM}} \\
	\textcolor{prussianblue}{\bfseries{Edificio}} /  \textcolor{trueblue}{\bfseries{Salón}}    & 	\textcolor{prussianblue}{\bfseries{FM9}} /  \textcolor{trueblue}{\bfseries{301}}   
\end{tabular} 
\end{mybox} \vspace{1cm}

\tableofcontents

\vspace{2cm}

\section{Calor específico de un cristal}



Analicemos \underline{desde el punto de vista clásico} el movimiento térmico de los átomos.

Consideremos la energía de cada átomo asumiéndolo como un oscilador armónico:


\begin{figure}[H]
    \centering
	    \includegraphics[width=0.4\textwidth]{1.png} % ajusta el ancho según convenga
    \caption{Descripción clara y concisa de la figura.}
    \label{fig:mi_figura} % opcional: para usar con \ref{fig:mi_figura}
\end{figure}

\begin{figure}[H]
    \centering
	    \includegraphics[width=0.7\textwidth]{2.png} % ajusta el ancho según convenga
    \caption{Descripción clara y concisa de la figura.}
    \label{fig:mi_figura} % opcional: para usar con \ref{fig:mi_figura}
\end{figure}

\begin{equation}
\begin{split}
\epsilon  &= \frac{1}{2} m \mathbf{v}^{2} + \frac{1}{2} \beta {\mathbf{\Delta {x}}}^{2}
\end{split}
\end{equation}

\begin{figure}[H]
    \centering
	    \includegraphics[width=0.7\textwidth]{3.png} % ajusta el ancho según convenga
    \caption{Descripción clara y concisa de la figura.}
    \label{fig:mi_figura} % opcional: para usar con \ref{fig:mi_figura}
\end{figure}

\textbf{Nueva bibliografía} Blackemore. Algunas cosas de clase no están en el Kittel

Estructura de \textbf{CS}. 

Por lo que la energía total  del sistema para un cristal de $N$ osciladores es:

\[ \boxed{  E=N\epsilon=3NK_{B}T }\]


Si ya tenemos la temperatura, si derivamos obtenemos:

\emph{El calor específico a volumen constante es}:


\[ C_{v} = \left( \frac{\partial E}{\partial T} \right)_{v} \] 



Y para el cristal es:

\[ \boxed{C_{v} = 3NK_{B}}  \]
 


Este es el resultado clásico establecido por la ley de \textbf{Dulong-Petit}

La ley de Dulong-Petit se estableció en base a experimentos. nos dice que el calor específico es constantes, es decir, no depende de $T$.

Alrededor de $T$ ambiente (300 K) y T$^{\prime}$s mayores, el calor específico de cristales es muy cercano al establecido por Dulong-Petit. Sin embargo, a finales del siglo \textbf{XIX} se volvió claro que a bajas T$^{\prime}$s su valor empieza a decaer. Hoy en día se sabe que tiende a cero a $T=0K$.


\begin{figure}[H]
    \centering
	    \includegraphics[width=0.7\textwidth]{4.png} % ajusta el ancho según convenga
    \caption{Descripción clara y concisa de la figura.}
    \label{fig:mi_figura} % opcional: para usar con \ref{fig:mi_figura}
\end{figure}


La suposición de osciladores armónicos clásicos no explica el comportamiento a bajas T$^{\prime}$s. Para tener mejores acuerdos con las onservaciones experimentales, en ese rango de temperaturas, se requiere un tratamiento cuántico.

\fbox{La idea es considerar osculadores cuánticos}

Para N osciladores independientes (en 3D):

\[ E= \sum\limits_{j}^{3N}\hbar \omega_{j} \left( n_{j}+\frac{1}{2} \right) = \sum\limits_{j}^{3N} \hbar\omega_{j}n_{j} + \sum\limits_{j}^{3N}  \frac{\hbar\omega_{j}}{2} = \sum\limits_{j}^{3N} \hbar\omega_{j}n_{j}+E_{0} \]

$E_{0} = \text{Energía de punto cero, es independiente de la temperatura}$

\underline{El modelo de Einstein}

En 1907 Einstein postuló la cuantización de la energía de las vibraciones de la red, es decir, una vibración de la red de energía $\nu$, sólo puee tomar valores múltiplos de $\hbar\omega$



\[ \langle \epsilon_{i} \rangle = \langle n_{i} \rangle \hbar\omega_{i} \]


\[\langle  n_{i} \rangle\] es el número promedio de partículas que se encuentra en el estado con energía $\langle \epsilon_{i} \rangle$

\begin{figure}[H]
    \centering
	    \includegraphics[width=0.7\textwidth]{5.png} % ajusta el ancho según convenga
    \caption{Descripción clara y concisa de la figura.}
    \label{fig:mi_figura} % opcional: para usar con \ref{fig:mi_figura}
\end{figure}


\begin{figure}[H]
    \centering
	    \includegraphics[width=0.7\textwidth]{6.png} % ajusta el ancho según convenga
    \caption{Descripción clara y concisa de la figura.}
    \label{fig:mi_figura} % opcional: para usar con \ref{fig:mi_figura}
\end{figure}


 Derivamos y obtenemos:
 
 El calor específico debido a las vibraciones de la red se ontiene derivando la expresión anterior:
 
 
 
 \[   C_{v} = \left( \frac{ \partial E }{\partial T}  \right)_{v} = 3NK_{B} \left( \frac{  \hbar \omega_{E} }{K_{B}T  } \right)^{2} \frac{ e^{  \hbar \omega_{E}/K_{B}T }   }{ (e^{ \hbar\omega_{E}/K_{B}T } -1  )^{2}  }  \]
 
 Definiendo:
 
 
  \[  x= \frac{  \hbar\omega }{K_{B}T} \]
 
 \[  x_{E} = \frac{  \hbar\omega_{E} }{K_{B}T} \]
 
 Y la función de Einstein como:
 
 \[F(x) = \frac{ x^{2}e^{x} }{ (e^{x}-1)^{2} } \]
 
 Tenemos
 
 \[ \boxed{   C_{v} = 3NK_{B}F(x_{E})  }  \]
 
 Hagamos un análisis de esta expresión a altas y bajas temperaturas.
 
 A temperaturas suficientemente altas, es decir cuando $K_{B}T >> \hbar\omega_{E}$, o bien $x_{E} = \hbar\omega_{E}/K_{B}T<<1$

Desarrollando en serie:


\[  F(x) = \frac{  x^{2}(1+x+x^{2}/2 + \ldots )  }{ (x+x^{2}/2+\ldots)^{2} } \approx 1  \]


Nos volamos los términos de mayor orden. 


Por lo que:

\[ \boxed{ C_{v} = 3NK_{B}F(x_{E})\approx 3NK_{B}  } \]

Es decir nos da el modelo de \textbf{Dulong-Petit} en altas temperaturas, empata bien \textcolor{Apple Green}{\textbf{\checkmark}}.

A temperaturas \textcolor{Cerulean}{\textbf{bajas}}, es decir, cuando $K_{B}<<\hbar\omega_{E}$, o bien $x_{E}=\hbar\omega_{E}/K_{B}T>>1$


En ese caso:


\[ F(x) = \frac{ x^{2} e^{x} }{ (e^{x}-1)^{2} } \approx x^{2}e^{-x}  \]

Y por ende:

\[ \boxed{ C_{v} = 3NK_{B}F(x_{E}) \approx 3NK_{B}x^{2}_{E}e^{-x_{E}}      }  \]


Tiene un decaimiento exponencial.

\textbf{NOTA:} 

\[  \frac{ \hbar\omega_{E}}{K_{B}} = \theta_{E} \] 

\begin{figure}[H]
    \centering
	    \includegraphics[width=0.7\textwidth]{7.png} % ajusta el ancho según convenga
    \caption{Descripción clara y concisa de la figura.}
    \label{fig:mi_figura} % opcional: para usar con \ref{fig:mi_figura}
\end{figure}


Es decir que a bajar temperaturas el modelo no empata. \textcolor{Cinnabar}{\xmark}

Investigar la \underline{temperatura de Einstein} para el diamante

\begin{figure}[H]
    \centering
	    \includegraphics[width=0.7\textwidth]{8.png} % ajusta el ancho según convenga
    \caption{Descripción clara y concisa de la figura.}
    \label{fig:mi_figura} % opcional: para usar con \ref{fig:mi_figura}
\end{figure}

\begin{figure}[H]
    \centering
	    \includegraphics[width=0.7\textwidth]{10.png} % ajusta el ancho según convenga
    \caption{Descripción clara y concisa de la figura.}
    \label{fig:mi_figura} % opcional: para usar con \ref{fig:mi_figura}
\end{figure}

\begin{figure}[H]
    \centering
	    \includegraphics[width=0.7\textwidth]{9.png} % ajusta el ancho según convenga
    \caption{Descripción clara y concisa de la figura.}
    \label{fig:mi_figura} % opcional: para usar con \ref{fig:mi_figura}
\end{figure}

\begin{figure}[H]
    \centering
	    \includegraphics[width=0.7\textwidth]{11.png} % ajusta el ancho según convenga
    \caption{Descripción clara y concisa de la figura.}
    \label{fig:mi_figura} % opcional: para usar con \ref{fig:mi_figura}
\end{figure}



Después del modelo de Debye empezamos con propiedades eléctricas de los materiales.



\end{document}